\section{\termBrokenCuspTrajectory{}模空间}
\label{sec:broken-cusp-trajectory-moduli-spaces}

设 $p,q\in L_0\cap L_1$。正如第~\ref{sec:splitting-ends} 节所指出的，要紧化
$\mathcal M(p,q;B)$，就需要添入由\personFloer{}流线与\termBubbleComponent{}之并构成的
对象；后者或者是伪全纯球面，或者是伪全纯圆盘。在这些\termIrreducibleComponent{}中，我们把由\personFloer{}流线构成的分支称为\emph{\termPrincipalComponent{}}。

\stfhIndex{I083}{Fuluoerjiandianliuxian}{弗洛尔尖点流线（Floer cusp-trajectory）}
先从只有唯一\termPrincipalComponent{}的情形开始。直观地说，类 $B$ 中一条从
$p$ 到 $q$ 的（未断裂）\emph{\personFloer{} \termCuspTrajectory{}}\footnote{译者注：原书本段已经
固定 $p,q\in L_0\cap L_1$，主分支也写成
$u_0\in\mathcal M(p,q;B_0)$，却把这条尖点流线说成从 $x$ 到 $y$。译文依上下文
统一为从 $p$ 到 $q$。}，是\personFloer{}流线
$u_0\in\mathcal M(p,q;B_0)$ 与有限个球面气泡和圆盘气泡之并；这些气泡附着于
$\mathbb R\times[0,1]$ 中有限多个点。\footnote{译者注：原书这里说附着的是
“有限多个球面气泡与圆盘气泡的模空间”，但下一页的纤维积以及定义 14.3.2
实际把稳定映射作为附着对象。译文据此写作气泡本身。}现在来精确定义这些对象。

先注意，模去定义~\ref{def:gamma-equivalent-fillings} 中的\personNovikov{}等价关系后，
$\pi_2(M)$ 以\termMonoidalAction{}作用于 $\pi_2(p,q)$，而对子
\[
  \bigl(\pi_2(M,L_0),\pi_2(M,L_1)\bigr)
\]
以显然的方式在 $\pi_2(p,q)$ 上给出\termBiMonoidalAction{}。以 $\alpha$ 表示
$\pi_2(M)$ 的元素，并对 $i=0,1$ 以 $\beta_i$ 表示 $\pi_2(M,L_i)$ 的元素。
上述作用在类 $B_0\in\pi_2(p,q)$ 上的像记为
\[
  B
  =
  B_0+\sum_\ell\alpha_\ell+\sum_j\beta_{0,j}+\sum_k\beta_{1,k}.
\]

\stfhIndex{I059}{Weixing}{位形（configuration）}
\begin{definition}
\label{def:strip-point-configuration}
$\mathbb R\times[0,1]$ 上的一个\termConfiguration{}，是由下列有限多个点组成的
集合：
\[
  \begin{array}{lll}
    (\tau_i,t_i)\in\mathbb R\times(0,1),
      & i=1,\ldots,m,\\
    (\tau_j,0)\in\mathbb R\times\{0\},
      & j=1,\ldots,n_0,\\
    (\tau_k,1)\in\mathbb R\times\{1\},
      & k=1,\ldots,n_1.
  \end{array}
\]
\stfhIndex{I020}{00-c-m-n0-n1}{$\mathfrak C_{(m;n_0,n_1)}$}
以 $\mathfrak C_{(m;n_0,n_1)}$ 表示所有这类位形的集合。
\end{definition}

一般把这样的三元组记为
\[
  C
  =
  \bigl(
    \{(\tau_i,t_i)\};
    \{(\tau_j,0)\},
    \{(\tau_k,1)\}
  \bigr).
\]
注意，只要 $C\neq\varnothing$，就不存在非平凡的全纯自同构。为方便起见，我们也把
空位形 $C=\varnothing$ 纳入讨论。称对子
\[
  (u,C)
  \in
  \widetilde{\mathcal M}(p,q;B)
  \times
  \mathfrak C_{(m;n_0,n_1)}
\]
\stfhIndex{I122}{DaibiaojiFuluoerliuxian}{带标记弗洛尔流线（marked Floer trajectory）}
为一条\termMarkedFloerTrajectory{}，其中
\[
  u:\mathbb R\times[0,1]\longrightarrow M,
  \qquad
  C\in\mathfrak C_{(m;n_0,n_1)}.
\]
在乘积
\[
  \widetilde{\mathcal M}_{(m;n_0,n_1)}(p,q;B)
  :=
  \widetilde{\mathcal M}(p,q;B)
  \times
  \mathfrak C_{(m;n_0,n_1)}
\]
上有自然的 $\mathbb R$-作用\footnote{译者注：原书在叙述这项作用时把乘积的
第一因子写成不带波浪号的 $\mathcal M(p,q;B)$，但作用公式、此前带标记流线的
定义以及紧随其后的商式都要求参数化空间
$\widetilde{\mathcal M}(p,q;B)$。原书又在公式 (14.3.40) 使用未曾定义的
$\widetilde{\mathcal M}_{(m;n_0,n_1)}$。译文把上述乘积明确记作这个带波浪号的
空间，并在下式取商。}
\[
  (u,C)\longmapsto(u\circ\tau_0,C\circ\tau_0),
  \qquad \tau_0\in\mathbb R,
\]
这里
\[
  (u\circ\tau_0)(\tau,t)=u(\tau-\tau_0,t),
\]
而 $C\circ\tau_0$ 定义为
\[
  \bigl(
    \{(\tau_i+\tau_0,t_i)\};
    \{(\tau_j+\tau_0,0)\},
    \{(\tau_k+\tau_0,1)\}
  \bigr).
\]
定义
\[
  \mathcal M_{(m;n_0,n_1)}(p,q;B)
  =
  \widetilde{\mathcal M}_{(m;n_0,n_1)}(p,q;B)/\mathbb R,
\]
并且当 $(m;n_0,n_1)=(0;0,0)$ 时定义
\[
  \mathcal M_\varnothing(p,q;B)
  =
  \widetilde{\mathcal M}(p,q;B)/\mathbb R.
\]
注意，当 $C\neq\varnothing$ 时，有自然的取值映射
\[
  \operatorname{ev}:
  \mathcal M_{(m;n_0,n_1)}(p,q;B)
  \longrightarrow
  M^m\times L_0^{n_0}\times L_1^{n_1},
\]
其定义为
\[
  \operatorname{ev}(u,C)
  =
  \bigl(
    \{u(\tau_i,t_i)\},
    \{u(\tau_j,0)\},
    \{u(\tau_k,1)\}
  \bigr).
\]
它与上述 $\mathbb R$-作用相容，因而定义良好。

与此前一样，以 $\overline{\mathcal M}_1(J_0;\alpha)$ 表示带一个标记点的亏格
$0$ 稳定映射，以 $\overline{\mathcal M}_1(L,J_0;\beta)$ 表示如第 9.5 节那样在
圆盘边界上带一个标记点的\termBorderedStableMap{}集合。考虑关于显然的取值映射所成的\termFiberProduct{}\footnote{译者注：原书公式 (14.3.40) 的三个乘积因子都把刚在
上文声明的 $\overline{\mathcal M}_1$ 漏写成 $\overline{\mathcal M}$；球面项又
写成 $J_{(\tau_i,t_i)}$，尽管本节固定的是 $J=\{J_t\}$；两个圆盘项则使用未指定的
$L$，而不是 $L_0,L_1$。定义 14.3.4 随即分别使用
$\overline{\mathcal M}_1(J_{t_i};\alpha_i)$、
$\overline{\mathcal M}_1(L_0,J_0;\beta_{0,j})$ 与
$\overline{\mathcal M}_1(L_1,J_1;\beta_{1,k})$，译文依该处补正。}
\begin{equation}
  \begin{aligned}
  &\widetilde{\mathcal M}
    \bigl(p,q;B_0;\{\alpha_i\},\{\beta_{0,j}\},\{\beta_{1,k}\}\bigr)
  \\
  &\quad:=
  \widetilde{\mathcal M}_{(m;n_0,n_1)}(p,q;B_0)_{\operatorname{ev}}
  \\
  &\qquad\mathbin{\times_{\operatorname{ev}}}
  \left(
    \begin{gathered}
      \displaystyle
      \prod_i\overline{\mathcal M}_1(J_{t_i};\alpha_i)
      \times
      \prod_j\overline{\mathcal M}_1(L_0,J_0;\beta_{0,j})
      \\
      {}\displaystyle\times
      \prod_k\overline{\mathcal M}_1(L_1,J_1;\beta_{1,k})
    \end{gathered}
  \right).
  \end{aligned}
\label{eq:cusp-trajectory-fiber-product}
\end{equation}
在
\[
  \widetilde{\mathcal M}
  \bigl(p,q;B_0;\{\alpha_i\},\{\beta_{0,j}\},\{\beta_{1,k}\}\bigr)
\]
上有 $\mathbb R$-作用；它同时平移附着在\termPrincipalComponent{}上的各个\termBubbleRoot{}。

\stfhIndex{I069}{Jiandianliuxian}{尖点流线（cusp-trajectory）}
\begin{definition}
\label{def:unbroken-cusp-trajectories}
定义从 $p$ 出发、终止于 $q$ 的\termCuspTrajectory{}集合为
\[
  \begin{aligned}
  &\mathcal M
  \bigl(p,q;B_0;\{\alpha_i\},\{\beta_{0,j}\},\{\beta_{1,k}\}\bigr)
  \\
  &\quad=
  \widetilde{\mathcal M}
  \bigl(p,q;B_0;\{\alpha_i\},\{\beta_{0,j}\},\{\beta_{1,k}\}\bigr)
  /\mathbb R.
  \end{aligned}
\]
进而定义
\[
  \mathcal{CM}(p,q;B)
  =
  \coprod
  \mathcal M
  \bigl(p,q;B_0;\{\alpha_i\},\{\beta_{0,j}\},\{\beta_{1,k}\}\bigr),
\]
其中不交并遍历所有满足
\[
  B=B_0+\sum_i\alpha_i+\sum_j\beta_{0,j}+\sum_k\beta_{1,k}
\]
的 $B_0,\{\alpha_i\},\{\beta_{0,j}\}$ 与 $\{\beta_{1,k}\}$，并在其上赋予\termStableMapTopology{}。
\end{definition}

把相应的映射分解记为
\[
  u
  =
  u^0
  \#
  \left(\coprod_i v_i\right)
  \#
  \left(\coprod_j w_{0,j}\right)
  \#
  \left(\coprod_k w_{1,k}\right).
\]
称 $u^0$ 为\termPrincipalComponent{}，称其余各分支为\termBubbleComponent{}，并称
这样的 $u$ 为一条\termCuspTrajectory{}。

\stfhSubIndex{I060}{Weixing}{位形（configuration）}{Duanlieweixing}{断裂位形（broken configuration）}
接着引入\termBrokenConfiguration{}的概念。对每个给定的整数 $n\in\mathbb N$，一般记
\[
  \underline n=\{1,\ldots,n\}.
\]

\begin{definition}
\label{def:broken-configuration}
一个\termBrokenConfiguration{}定义为对子
\[
  (\Theta,o),
\]
其中 $\Theta$ 是有限多个
$\Theta_\ell\cong\mathbb R\times[0,1]$ 的连通并，而
\[
  o=(o^{\mathrm{sphere}};o^{\mathrm{disc},0},o^{\mathrm{disc},1}),
\]
且
\[
  o^{\mathrm{sphere}}:
  \underline m
  \longrightarrow
  \bigcup_{\ell=1}^a\operatorname{Int}\Theta_\ell
\]
是单射，
\[
  o^{\mathrm{disc},0}:
  \underline{n_0}
  \longrightarrow
  \bigcup_{\ell=1}^a\partial_0\Theta_\ell,
  \qquad
  o^{\mathrm{disc},1}:
  \underline{n_1}
  \longrightarrow
  \bigcup_{\ell=1}^a\partial_1\Theta_\ell
\]
也是单射。这里对每个 $\ell=1,\ldots,a$，在所选标准坐标下
记\footnote{译者注：原书先把各分支写成
$\Theta_\ell\cong\mathbb R\times[0,1]$，这里却只写
$\Theta_a=\mathbb R\times[0,1]$，同时仍以 $\ell$ 标记边界。译文依此前的分支族把
下标统一为 $\ell$，并把等号理解为选定标准坐标后的识别。}
\[
  \Theta_\ell=\mathbb R\times[0,1],
  \qquad
  \partial_i\Theta_\ell=\mathbb R\times\{i\},
  \quad i=0,1.
\]
常记
\[
  o^{\mathrm{sphere}}(i)=(\tau_i,t_i),
  \qquad
  o^{\mathrm{disc},0}(j)=(\tau_j,0),
  \qquad
  o^{\mathrm{disc},1}(k)=(\tau_k,1),
\]
并把 $o$ 与其像点等同起来。
\end{definition}

\stfhIndex{I180}{WendingduanlieFuluoerliuxian}{稳定断裂弗洛尔流线（stable broken Floer trajectory）}
\begin{definition}
\label{def:stable-broken-floer-trajectory}
设 $J=\{J_t\}_{0\leq t\leq1}$，并设
$p,q\in L_0\cap L_1$。\footnote{译者注：原书此前和定义正文都使用 $p,q$，
定义首句却写成 $x,y,\in L_0\cap L_1$，其中还多出一个逗号。译文统一为
$p,q$。}从 $p$ 到 $q$ 的一条\termStableBrokenFloerTrajectory{}，是满足下列条件的
三元组\footnote{译者注：原书列出的四项条件没有要求气泡根处的取值与主分支取值
相等，也没有给出总类 $B$ 与各主分支、球面及圆盘类之和的关系；此外
$u_\ell\in\mathcal M(x_\ell,x_{\ell+1})$ 是平移商类，而 $o$ 使用各条带上的绝对
$\tau$-坐标。因而这些条件本身不足以确定后文
$\overline{\mathcal M}(p,q;B)$ 所需的附着关系、同伦类和参数化数据。原书没有给出
可据以唯一补全的额外条件，译文保留这一缺口。}
\[
  u
  =
  \bigl(
    (u_1,\ldots,u_a);
    (\sigma_1,\ldots,\sigma_m),
    (\gamma_1^0,\ldots,\gamma_{n_0}^0),
    (\gamma_1^1,\ldots,\gamma_{n_1}^1);
    o
  \bigr).
\]
\begin{enumerate}
  \item 存在
  $x_1,\ldots,x_{a+1}\in L_0\cap L_1$，其中
  $x_1=p$、$x_{a+1}=q$\footnote{译者注：原书在公式 (14.3.41) 中仅给 $p$ 加帽，
  写作 $u_1(-\infty)=\widehat p$，而本节的 $p,q$ 都是 $L_0\cap L_1$ 中的交点，
  且正端仍写 $q$。译文删去这个孤立而未定义的帽号。}，使得对
  $\ell=1,\ldots,a$，
  \[
    u_\ell\in\mathcal M(x_\ell,x_{\ell+1}),
  \]
  并且\footnote{译者注：原书只令
  $i=1,\ldots,a-1$ 时 $u_i\in\mathcal M(x_i,x_{i+1})$，从而漏掉 $u_a$；
  这些中间交点 $x_i$ 也没有量化。译文依端点式和下一行的相邻匹配唯一补全。}
  \begin{equation}
    \begin{gathered}
      u_1(-\infty)=p,
      \qquad
      u_a(\infty)=q,\\
      u_\ell(\infty)=u_{\ell+1}(-\infty),
      \qquad \ell=1,\ldots,a-1.
    \end{gathered}
  \label{eq:broken-trajectory-matching}
  \end{equation}
\stfhIndex{I098}{Kezhanhe}{可粘合（gluable）}
  称~\eqref{eq:broken-trajectory-matching} 为\termMatchingCondition{}；若一对\personFloer{}流线 $(u,u')$ 满足该条件，就称它们\termGluable{}。

  \item 对 $i=1,\ldots,m$，
  \[
    \sigma_i\in\overline{\mathcal M}_1(J_{t_i};\alpha_i).
  \]

  \item 对 $j=1,\ldots,n_0$ 以及 $k=1,\ldots,n_1$，
  \[
    \gamma_j^0
    \in\overline{\mathcal M}_1(L_0,J_0;\beta_j^0),
    \qquad
    \gamma_k^1
    \in\overline{\mathcal M}_1(L_1,J_1;\beta_k^1).
  \]

  \item 对每个 $\ell=1,\ldots,a$，或者映射 $u_\ell$ 是\termNonStationary{}的，
  或者
  \[
    \Theta_\ell\cap\operatorname{Im}o\neq\varnothing.
  \]
\end{enumerate}
\end{definition}

把 $u$ 的定义域简记为 $\Theta_u$；它由一个\termBrokenConfiguration{}与稳定映射 $\sigma_i,\gamma_j^0,\gamma_k^1$ 的定义域，
在给定标记点处粘到 $o$ 的像点而成。\footnote{译者注：原书称
$\Theta_u$ 是断裂位形与各稳定映射定义域的“乘积”，紧接着却明确说
这些定义域粘到标记点且整体是一个并。译文依后一描述写作粘合而成。}它是下列对象的并：
有限多个
\[
  \mathbb R\times[0,1]
\]
的副本，即\termPrincipalComponent{}，以及闭或带边亏格 $0$ \personRiemann{}曲面的\termSemiStableCurve{}，即\termBubbleComponent{}；后者的\termBubbleRoot{}按照\termDistributionMap{} $o$ 附着于 $\Theta_u$ 的主分支上。

现在记三元组与相应的标号集为
\[
  I=(m,n_0,n_1),
  \qquad
  \underline I
  =
  \underline m\coprod\underline{n_0}\coprod\underline{n_1},
\]
并以
\[
  \rho:\underline I\longrightarrow\underline I
\]
表示保持\termCoproduct{}
$\underline m\coprod\underline{n_0}\coprod\underline{n_1}$ 各因子的双射。
\footnote{译者注：原书称 $I$ 是“这些集合中的任一个”，随后却同时以
$I=(m,n_0,n_1)$ 表示三元组、以带下划线的 $I$ 表示三个标号集的余积。译文把二者
分开陈述。}

若存在这样的双射 $\rho:\underline I\to\underline I$，使得
\[
  u_i'=u_{\rho(i)},
  \qquad
  \sigma_j'=\sigma_{\rho(j)},
  \qquad
  \gamma_k'=\gamma_{\rho(k)},
\]
就把两条稳定断裂流线 $u$ 与 $u'$ 等同；仍滥用记号，以
$u$ 表示相应的等价类。\footnote{译者注：这里保留原书的三个等式。
但 $\rho$ 的定义域只标记 $m$ 个球面气泡和两族圆盘气泡，并不标记 $a$ 个主分支，
故首式的 $u_{\rho(i)}$ 类型不成立；次式又以圆盘标号 $j$ 索引球面族
$\sigma$，末式没有区分 $\gamma^0,\gamma^1$。原书也未说明分布映射 $o$ 如何随
$\rho$ 变化。因此所印条件不能完整定义所声称的等价关系，而仅据本段无法唯一重构。}

最后，定义
\[
  \overline{\mathcal M}(p,q;B)
\]
为这些同构类的集合。与此前研究稳定映射模空间时一样，在
$\overline{\mathcal M}(p,q;B)$ 上定义一个拓扑，使它成为紧\personHausdorff{}空间；该
拓扑在
\[
  \mathcal M(p,q;B)
  =
  \widetilde{\mathcal M}(p,q;B)/\mathbb R
\]
上诱导的拓扑，是 $\widetilde{\mathcal M}(p,q;B)$ 的\termStrongCinfinityTopology{}所给出的商拓扑。

这里要强调，这个紧化对于横截拉格朗日子流形对 $L_0,L_1$ 以及
$J=\{J_t\}_{0\leq t\leq1}$ 的\emph{任意}选择，都是作为拓扑空间定义的。即使对一般的
对子 $(L_0,L_1)$ 通有地选择 $J$，拓扑空间
$\overline{\mathcal M}(p,q;B)$ 在通常意义下也不是光滑流形，而只能在\termKuranishiStructure{}的意义下理解 (FOn99)。不过，有一类特殊的拉格朗日子流形，
即\emph{\termMonotoneLagrangianSubmanifold{}}；对它们而言，当
\[
  J
  \in
  \mathfrak j^{\mathrm{reg}}(L_0,L_1)
  \subset
  \mathfrak j_\omega
\]
取自一个稠密集合时，$\overline{\mathcal M}(p,q;B)$ 成为具有预期维数的
光滑\termOrbifoldWithCorners{}。
